% 08_compare.tex - compare the six runs
\section{Compare the runs}

\step{Compare the six runs}
The script \cmd{compare\_runs.py} prints the mean values of each run and draws the recharge that reaches the water table,
the water-table depth of the first HRU given, and the streamflow:
\term{t3_compare.txt}
\fig{t3_compare.png}{\cmd{compare\_runs.png}: the six runs. Recharge reaching the water table (30-day means), water-table
depth under the north slope of the lower subbasin (HRU~7), and streamflow at the outlet.}
\why the table shows three effects.
\begin{itemize}
\item \textbf{From CONCEPTUAL to PHYSICAL.} In the conceptual run, \V{t3:conceptual:leaves}~m\textsuperscript{3}/d leave
  the subsoil and all of it reaches the aquifer: \V{t3:conceptual:reaches}~m\textsuperscript{3}/d the water table, the rest
  the pond's specified-head cells, which pass it straight to the pond (Workshop~2, Step~13);
  \V{t3:conceptual:seep}~m\textsuperscript{3}/d seep back out. With the soil
  zone, \V{t3:physical:leaves}~m\textsuperscript{3}/d leave the subsoil, because on the valley floor the soil now drains
  into an aquifer whose water table is in the soil. But \V{t3:physical:rejected}~m\textsuperscript{3}/d of it is
  rejected, and \V{t3:physical:aqsoil}~m\textsuperscript{3}/d rise back from the aquifer into the soil. A two-way exchange
  replaces the one-way seepage drain: the seepage to Raven falls to \V{t3:physical:seep}~m\textsuperscript{3}/d, from the
  lake HRU only. The water table under the valley floor (HRU~6) lies at \V{t3:physical:wt6}~m instead of
  \V{t3:conceptual:wt6}~m, and more water reaches the creek through the soil (\V{t3:physical:q} instead of
  \V{t3:conceptual:q}~m\textsuperscript{3}/s).
\item \textbf{The unsaturated columns hold water.} At the end of the run, RICHARDS stores \V{t3:richards:uz}~million
  m\textsuperscript{3} above its equilibrium profile (\V{t3:uzmm}~mm over the \V{t3:slopearea}~km\textsuperscript{2} of
  slopes). LAG stores only \V{t3:lag:uz}~million m\textsuperscript{3} above the same profile, and KINEMATIC
  \V{t3:kinematic:uz}~million m\textsuperscript{3} (counted from the residual water content). The difference is physical.
  The columns start at rest, in equilibrium with a deep water table, and RICHARDS must wet this dry profile before it can
  carry the recharge. Over the three years it holds back \V{t3:withheld}\% of what enters the columns. So only
  \V{t3:richards:reaches}~m\textsuperscript{3}/d reach the water table, against \V{t3:kinematic:reaches} with KINEMATIC,
  and the water table under the slopes lies on average \V{t3:wtdrop}~m deeper than without columns. Most of this is spin-up (see the
  note below). The recharge also arrives later and smoother (Figure~\ref{fig:t3vz}). The flow out of the columns follows
  the flow into them best \V{t3:richards:lag}~days later with RICHARDS, \V{t3:kinematic:lag} days with KINEMATIC and
  \V{t3:lag:lag} days with LAG.
\item \textbf{Depth classes change little here}: \V{t3:classes:reaches} against \V{t3:richards:reaches}
  m\textsuperscript{3}/d reach the water table, and the run takes longer. Depth classes matter where the unsaturated
  thickness within a cell varies a lot compared with its mean, as on uneven ground over a shallow water table.
\end{itemize}
The script \cmd{plot\_vadose.py} draws the unsaturated columns of the three methods:
\term{t3_plot_vadose.txt}
\fig{t3_vadose.png}{\cmd{vadose.png}: the unsaturated columns of the three methods: what enters them and what reaches the water table
(30-day means), and the water they store.\label{fig:t3vz}}
\begin{note}[A deep unsaturated zone needs a spin-up]
A \emph{spin-up} is a run that brings the model's stores to a realistic state before the period you study. The columns
start at rest (in equilibrium with the water table) and fill during the run. With RICHARDS, the water table under the
slopes falls from \V{t3:physical:wt7}~m (no columns) to a mean of \V{t3:richards:wt7}~m, and the storage is still growing
after three years. Look at the column \cmd{unsaturated zone storage} of \cmd{GWBudget.csv}: as long as it trends up or
down, the run is still spinning up. There are two remedies. Run long enough and start the real run from its
\cmd{solution.rvc} (Workshop~5 shows how; the columns are saved there). Or start each column at its steady profile for a
typical recharge with \cmd{:SoilZoneVadoseInitialFlux} (in mm/d).
\end{note}
